\section[Fargues--Fontaine 曲线 II]{Fargues--Fontaine 曲线 II（11 月 5 日）}
\srcpage{12}

\subsection{回顾}

上次我们考虑了等特征局部域 $E=\mathbb{F}_q((t))$ 的 Fargues--Fontaine 曲线。

设 $C/\mathbb{F}_q$ 是完备代数闭非阿基米德域，我们考虑了
\[
\Spa E\times_{\Spa\mathbb{F}_q}\Spa C=\mathbb{D}^*_C,
\]
即 $C$ 上的去心开单位圆盘。它有 $\varphi_C$ 的作用，而我们定义了
$X_{C,E}:=\mathbb{D}^*_C/\varphi_C^{\mathbb{Z}}$。

\subsection{经典点}

有范畴等价
\[
\{\text{刚性解析簇}/C\}\cong\{\text{「局部有限型」adic 空间}/\Spa C\}.
\]
该等价由 $X(C)\leftarrow X$ 给出，其中 $X(C)\subset|X|$ 是「经典点」。
局部地 $X=\Spa A$，其中 $A=C\langle T_1,\dots,T_n\rangle/I$，并且
\[
X(C)=\{(x_1,\dots,x_n)\in C^n: |x_i|\le1,\ f(x_1,\dots,x_n)=0\}=:|\mathrm{Sp}\,A|.
\]

Tate 定义了 Grothendieck 拓扑 $(\mathrm{Sp}\,A)_{\mathrm{rig}}$，它由「拟紧容许开集」组成，
等价于 $\Spa A$ 上由拟紧开集给出的拓扑。用这些术语，Tate 的
$(\mathrm{Sp}\,A)_{\mathrm{rig}}$ 的「容许覆盖」与 $\Spa A$ 的覆盖相对应。

特别地，$\Spa A$ 的拟紧开子集完全由它们的经典点决定。

对 $\mathbb{D}^*_C$，经典点是 $\{x\in C: 0<|x|<1\}$。Frobenius $\varphi_C$ 在经典点集上的作用为
$x\mapsto x^{1/q}$。（它不是 $x\mapsto x^q$，原因在于 $\varphi_C$ 是通过它在 $C$ 上的作用而起作用的。）

对任何连通的 affinoid $\Spa A\subset\mathbb{D}^*_C$，$A$ 是主理想整环。这反映了 $\mathbb{D}^*_C$ 的一维性。
可以手工验证：对应于 $x$ 的极大理想由 $T-x$ 生成。

通过下降，我们也可以定义 $X_{C,E}$ 的经典点。它们是
\[
X^{\mathrm{cl}}_{C,E}=\{x\in C: 0<|x|<1\}/\varphi_C .
\]
同样，$X_{C,E}$ 的任何连通 affinoid 子集都是某个主理想整环的 adic 谱。

\subsection{混合特征的 Fargues--Fontaine 曲线}

现在我们转向 $E$ 是 $\mathbb{Q}_p$ 的有限扩张的情形。仍设 $C/\mathbb{F}_q$ 如前
（完备代数闭非阿基米德域）。

什么是「$\Spa E\times_{\Spa\mathbb{F}_q}\Spa C$」？当然，直接的问题是：由于 $E$ 的特征为 $0$，
不存在映射 $\Spa E\to\Spa\mathbb{F}_q$。

想法是：在特征 $p$ 中，我们通过取 $R[[t]]$ 把任意 $\mathbb{F}_q$-代数 $R$ 形变到 $\mathbb{F}_q[[t]]$；
现在我们要把 $R$ 形变到混合特征。

注意：若 $R$ 是完美 $\mathbb{F}_q$-代数，则存在唯一的（至多差唯一同构）提升
$\tilde R/\mathcal{O}_E$，它平坦、$\pi$-adic 完备，并且 $\tilde R/\pi\xrightarrow{\sim}R$。
一种构造方式是借助（$p$-典型、分歧的）Witt 向量
\[
\tilde R=W_{\mathcal{O}_E}(R):=W(R)\otimes_{W(\mathbb{F}_q)}\mathcal{O}_E .
\]

Witt 向量的构造背后的想法如下。由于假设 $R$ 完美，$R$ 中每个元素都有 $p$ 次幂根。
\srcpage{13}由于取 $p$ 次幂会改善同余性，可以对 $p$ 次幂根的任意提升取指数，
从而写出典范的提升。更确切地说，存在 Teichmüller 映射
\[
[\cdot]:R\to\tilde R=W_{\mathcal{O}_E}(R)
\]
把 $x$ 送到 $\lim_{n\to\infty}\tilde x_n^{p^n}$，其中 $\tilde x_n\in\tilde R$ 是 $x^{1/p^n}$ 的任一提升。
它是乘性的，但当然不是加性的。

$\tilde R$ 的每个元素都有唯一的表达式
\[
\sum_{n\ge0}[r_n]\pi^n,\qquad r_n\in R .
\]
直观上，$\tilde R$ 的元素看起来像「$R$ 上以 $\pi$ 为变量的幂级数」。

在混合特征情形中，$\Spa\mathbb{F}_q[[t]]\times_{\Spa\mathbb{F}_q}\Spa\mathcal{O}_C
=\Spa\mathcal{O}_C[[t]]$ 的类比物应当是
\[
\text{「}\Spa\mathcal{O}_E\times_{\Spa\mathbb{F}_q}\Spa\mathcal{O}_C\text{」}
:=\Spa W_{\mathcal{O}_E}(\mathcal{O}_C).
\]
在等特征中我们实际处理的是 $\Spa\mathbb{F}_q((t))\times_{\Spa\mathbb{F}_q}\Spa C$。
它由要求 $t\ne0$ 与 $\varpi\ne0$ 得到。所以 $\Spa\mathbb{F}_q((t))\times_{\Spa\mathbb{F}_q}\Spa C$
在混合特征中的类比物是
\[
\text{「}\Spa E\times_{\Spa\mathbb{F}_q}\Spa C\text{」}:=Y_{C,E}
:=\{\pi\ne0,\ [\varpi]\ne0\}\subset\Spa W_{\mathcal{O}_E}(\mathcal{O}_C).
\]
它仍带有 $\varphi_C$ 的作用，该作用自由且完全间断。（不再有来自 $E$ 的 Frobenius 作用。）

\begin{definition}
\emph{Fargues--Fontaine 曲线}是 $X_{C,E}:=Y_{C,E}/\varphi_C^{\mathbb{Z}}$。
\end{definition}

\begin{remark}
$X_{C,E}$ 的几何在某种意义上与等特征情形极其相似，下面我们将看到。
一个区别是 $Y_{C,E}$ 远不那么显式。特别地，$Y_{C,E}$ 生活在 $\Spa E$ 上，
因而没有到 $\Spa C$ 的结构映射。
\end{remark}

\begin{theorem}[Fargues--Fontaine, Kedlaya]
\begin{enumerate}[label=(\arabic*), leftmargin=2.6em]
  \item 存在经典点的概念 $Y^{\mathrm{cl}}_{C,E}\subset Y_{C,E}$，使得对任何连通 affinoid 开子集
        $\Spa A\subset Y_{C,E}$，$A$ 是主理想整环，且
        $\mathrm{Spm}(A)\xrightarrow{\sim}\Spa A\cap Y^{\mathrm{cl}}_{C,E}$。
  \item 对任何经典点 $y\in Y^{\mathrm{cl}}_{C,E}$，存在 $x\in C$ 满足 $0<|x|<1$，使得
        $y=V(\pi-[x])$。（但要注意 $x$ 并不唯一。）\footnote{若 $C$ 不是代数闭的，
        我们就不能把元素取成 $\pi-[x]$ 这种简单形式。}
  \item 对任何经典点 $y\in Y^{\mathrm{cl}}_{C,E}$，$y$ 处的完备剩余域\footnote{%
        若 $(A,A^+)$ 是 Huber 对且 $x\in\Spa(A,A^+)$，则得到赋值 $|\cdot|_x:A\to\Gamma_x\cup\{0\}$。
        其核 $\mathfrak{p}_x:=\{f\in A: |f|_x=0\}\subset A$ 是素理想。于是 $x$ 处的\emph{完备剩余域}
        是 $\kappa_x:=\mathrm{Frac}(A/\mathfrak{p}_x)$ 关于 $x$ 的完备化。
        在我们当前的情形中，既不需要取完备化，也不需要取分式域。}
        （由 (1)，它就是 $A$ 的完备剩余域，其中 $\Spa A$ 是 $y$ 的 affinoid 邻域）
        是一个完备代数闭域 $C(y)$，带有典范同构
        \[
        \underbrace{C(y)^\flat}_{\text{「$C(y)$ 的倾斜」}}\xrightarrow{\sim}C .
        \]
        这诱导出双射
        \[
        Y^{\mathrm{cl}}_{C,E}\xrightarrow{\sim}
        \{C^\#/E:\ C^\#\text{ 是 }C\text{ 的 untilt}\}.\footnote{按定义，$C$ 的一个「untilt」
        包含其倾斜与 $C$ 之间的同构这一数据。}
        \]
\end{enumerate}
\end{theorem}

\srcpage{14}

\begin{remark}
回忆在等特征中，$Y_{C,E}$ 的经典点与 $\{x\in C: 0<|x|<1\}$ 显式地一一对应。
所以 (2) 与它的等特征类比略有不同。(3) 的类比在等特征中成立，但它是空洞的，
因为在那里倾斜是恒等函子。
\end{remark}

\subsection{倾斜}

对满足 $|p|_K<1$ 的完备代数闭非阿基米德域 $K$，可以定义一个完备代数闭非阿基米德（特征 $p$）域
$K^\flat$。作为乘性拓扑幺半群，$K^\flat\cong\varprojlim_{x\mapsto x^p}K$。

在这个逆极限中投影到第一个坐标，诱导出乘性（但显然非加性）的映射
$K^\flat\to K$ 与 $\mathcal{O}_{K^\flat}\to\mathcal{O}_K$，记作 $x\mapsto x^\#$。

为看出环结构，注意\footnote{它是同构这一点，归根结底依赖于构造 Teichmüller 映射时所用的同样的初等考量。}
\[
\mathcal{O}_{K^\flat}:=\varprojlim_{x\mapsto x^p}\mathcal{O}_K
\xrightarrow{\sim}\varprojlim_{x\mapsto x^p}\mathcal{O}_K/p .
\]
现在右边有显然的环结构。

\subsection{定理 3.3 的证明}

\subsubsection{第 1 步}

我们先构造单射
\[
\{C^\#/E:\ C^\#\text{ 是 }C\text{ 的 untilt}\}\longrightarrow|Y_{C,E}|.
\]
设 $C^\#$ 是 $C$ 的一个 untilt。投影
\[
\mathcal{O}_C\cong\varprojlim_{x\mapsto x^q}\mathcal{O}_{C^\#}
\xrightarrow{\mathrm{pr}_1}\mathcal{O}_{C^\#}
\]
把 $x$ 送到 $x^\#$，它诱导映射\footnote{我们将称之为「Fontaine 映射」。}
$\theta:W(\mathcal{O}_C)\to\mathcal{O}_{C^\#}$，把 $\sum_{n\ge0}[x_n]\pi^n$ 送到
$\sum_{n\ge0}x_n^\#\pi^n$。

于是我们得到 $\Spa\mathcal{O}_{C^\#}\to\Spa W_{\mathcal{O}_E}(\mathcal{O}_C)$。
$\theta$ 的满射性说明这是单射。与概形的情形一样，$\Spa\mathcal{O}_{C^\#}$ 有两个点——特殊点与一般点——
而这个映射诱导 $\Spa C^\#\to Y_{C,E}\subset\Spa W_{\mathcal{O}_E}(\mathcal{O}_C)$。
其像是 $Y_{C,E}$ 中的一个点 $y$，该点处的完备剩余域是 $C^\#$。

这就给出想要的映射
\[
\{C^\#/E:\ C^\#\text{ 是 }C\text{ 的 untilt}\}\hookrightarrow|Y_{C,E}|.
\]
为得到单射性，需要追踪一遍以看出 untilt 可以由 $\theta$ 重构。
定义 $Y^{\mathrm{cl}}_{C,E}$ 为这个像——这样我们就已经证明了定理的 (3)。

\subsubsection{$Y_{C,E}$ 的倾斜}

设 $E_\infty=E(\pi^{1/p^\infty})^\wedge:=(\bigcup_n E(\pi^{1/p^n}))^\wedge$。
这是一个「完满域（perfectoid field）」，意思是在 $\mathcal{O}_{E_\infty}/p$ 上 $x\mapsto x^p$ 是满射。

则倾斜是 $E^\flat_\infty\cong\mathbb{F}_q((t^{1/p^\infty}))$。这个等同通过令
$t:=(\pi,\pi^{1/p},\dots)\in\varprojlim_{x\mapsto x^p}E_\infty=E^\flat_\infty$ 来实现。

\begin{lemma}
完满空间 $(Y_{C,E}\times_{\Spa E}\Spa E_\infty)^\flat$ 典范同构于
$\mathbb{D}^*_C\times_{\Spa\mathbb{F}_q((t))}\Spa\mathbb{F}_q((t^{1/p^\infty}))$。
此外，经典点在这个对应下双射。
\end{lemma}

\srcpage{15}

注意，由于取完美化不影响底层拓扑空间，我们有
\[
|\mathbb{D}^*_C|\xleftarrow{\sim}
|\mathbb{D}^*_C\times_{\Spa\mathbb{F}_q((t))}\Spa\mathbb{F}_q((t^{1/p^\infty}))| .
\]
又倾斜是同胚，故
\[
|\mathbb{D}^*_C\times_{\Spa\mathbb{F}_q((t))}\Spa\mathbb{F}_q((t^{1/p^\infty}))|
\cong|Y_{C,E}\times_{\Spa E}\Spa E_\infty| .
\]
后者当然仍容许一个到 $|Y_{C,E}|$ 的映射，容易验证它是满射，尽管它不是同胚。
\[
\begin{tikzcd}[column sep=large, row sep=large]
{|Y_{C,E}\times_{\Spa E}\Spa E_\infty|} \arrow[r, equal] \arrow[d]
  & {|\mathbb{D}^*_C\times_{\Spa\mathbb{F}_q((t))}\Spa\mathbb{F}_q((t^{1/p^\infty}))|}
    \arrow[d, "\sim"] \\
{|Y_{C,E}|} & {|\mathbb{D}^*_C|}
\end{tikzcd}
\]
复合起来就得到 $\{x\in C: 0<|x|<1\}=\mathbb{D}^{*,\mathrm{cl}}_C\twoheadrightarrow Y^{\mathrm{cl}}_{C,E}$，
而追一遍定义便容易看出它把 $x$ 送到 $V(\pi-[x])$。

\subsubsection{关于完满空间的插叙}

\begin{definition}
一个\emph{完满 Tate 环}是一个完备 Tate 环 $A$（即存在拓扑幂零单位 $\varpi\in A$
以及定义环 $A_0$），使得
\begin{itemize}
  \item 存在 $\varpi$ 在 $A^\circ$ 中满足 $\varpi^p\mid p$，且
  \item $A^\circ$ 是 $\varpi$-adic 的（等价地，$A^\circ$ 是定义环），且
  \item $x\mapsto x^p$ 在 $A^\circ/p$ 上是满射。
\end{itemize}
一个\emph{完满空间}（perfectoid space）是被形如 $\Spa(A,A^+)$ 的开集覆盖的 adic 空间，
其中 $A$ 是完满 Tate 环。
\end{definition}

\begin{example}
完满域包括 $E_\infty$、$C$、$\mathbb{F}_q((t^{1/p^\infty}))$。
一个不是域的完满环的例子是 $A=C\langle T^{1/p^\infty}\rangle$。
\end{example}

\begin{example}
若 $A/\mathbb{F}_p$ 是特征 $p$ 的 Tate 环，则 $A$ 是完满的当且仅当 $A$ 是完美的。
\end{example}

倾斜可以推广到完满环\footnote{对我们来说，「完满环」总是指「完满 Tate 环」。
另有一个不同的概念叫作「整完满环（integral perfectoid ring）」，它在本课程中不起作用。}。
若 $A$ 完满，则倾斜为 $A^\flat=\varprojlim_{x\mapsto x^p}A$（配以适当的加法）。
注意存在乘性幺半群映射 $A^\flat\xrightarrow{\mathrm{pr}_1}A$，记作 $a\mapsto a^\#$。

\begin{example}
$E_\infty\langle T^{1/p^\infty}\rangle^\flat=E^\flat_\infty\langle T^{1/p^\infty}\rangle$。
\end{example}

倾斜也推广到完满空间 $X\mapsto X^\flat$，通过粘合 $\Spa(A,A^+)\mapsto\Spa(A^\flat,A^{\flat+})$。

\begin{theorem}[倾斜等价]
\begin{enumerate}[label=(\arabic*), leftmargin=2.6em]
  \item 存在同胚 $|X|\xrightarrow{\sim}|X^\flat|$，把 $x$ 送到 $x^\flat$，其中赋值由
        $|f(x^\flat)|:=|f^\#(x)|$ 给出。口号：「倾斜保持拓扑空间」。
  \item 给定完满空间 $X$，倾斜诱导范畴等价
        \[
        \{\text{完满空间 }Y/X\}\xrightarrow{\sim}\{\text{完满空间 }Y'/X^\flat\}.
        \]
        特别地，在固定一个基完满空间之后，untilt 是典范的。
  \item 若 $X=\Spa(A,A^+)$ 且 $X^\flat=\Spa(A^\flat,A^{\flat+})$，则 $X$ 与 $X^\flat$ 的
        Zariski 闭子集（即某个理想的零点轨迹）在倾斜下相对应：
        \[
        \{Z\subseteq|X|\}\leftrightarrow\{Z^\flat\subseteq|X^\flat|\}.
        \]
\end{enumerate}
\end{theorem}

\srcpage{16}

\begin{example}
挑战：设 $X=\Spa C^\#\langle T^{1/p^\infty}\rangle\supseteq Z=V(T-1)$。证明 $Z^\flat$ 是 Zariski 闭的。
\end{example}

\subsubsection{回到证明}

为证明引理 3.5，手工算出\footnote{这是整完满环的一个例子；对整完满环，倾斜也表现良好，
尽管它不是 Tate 的。}
\[
(W_{\mathcal{O}_E}(\mathcal{O}_C)\widehat{\otimes}_{\mathcal{O}_E}\mathcal{O}_{E_\infty})^\flat
\cong\mathcal{O}_C[[t^{1/p^\infty}]] .
\]
基本上我们需要把左边模 $p$ 与右边模 $t$ 等同起来。为此注意
\[
(W_{\mathcal{O}_E}(\mathcal{O}_C)\widehat{\otimes}_{\mathcal{O}_E}\mathcal{O}_{E_\infty})^\flat/p
\cong\mathcal{O}_C[[p^{1/p^\infty}]]^\wedge/p ,
\]
它同构于 $\mathcal{O}_C[[t^{1/p^\infty}]]/t$。

由于经典点是通过倾斜定义的，不难看出
$W_{\mathcal{O}_E}(\mathcal{O}_C)\widehat{\otimes}_{\mathcal{O}_E}\mathcal{O}_{E_\infty}$
与 $\mathcal{O}_C[[t^{1/p^\infty}]]$ 的经典点相对应，并且它们的理想有非常显式的生成元。
这就建立了定理 3.3 的 (2)。

剩下要建立 (1)，即主理想性。基本上我们要证明的是：$A$ 上的任何函数只能在经典点处消失。
这再次通过倾斜来证明。在倾斜后的图景中这是清楚的，因为 Zariski 闭子集是什么已经清楚了。
这说明 $\mathrm{Spm}(A)\hookrightarrow\Spa A\cap Y^{\mathrm{cl}}_{C,E}$。
但我们在 (2) 中已经得到了满射性。由于各点的理想是主理想，容易推出每个理想都是主理想。
